\providecommand{\anE}{\varepsilon}
\providecommand{\anD}{\mathcal D}
\providecommand{\anH}{\widehat}

\section{Continuous completion as a reusable interface}
\label{sec:continuous-completion}

The present section isolates three operations that turn finite analytic
estimates into continuous mathematics: extending a map to a completion,
gluing local maps with explicit weights, and identifying a completed point
from an approximate solution procedure.  The distinction between these
operations matters.  A Cauchy estimate by itself does not identify the
completed value of a function, and a collection of correct local answers
does not by itself give a continuous global answer.

There is already substantial relevant work in the inspected repository.
\path{CoherentMinimizer.cauchy} derives a regular name from a quantitative
uniqueness modulus; its source explicitly leaves canonical evaluation and
uniqueness of the completed point outside that theorem.
\path{StrongMinimizer.selected_pair_stable} obtains a distance estimate
for actual finite-net minimizers from a strong-convexity certificate.
\path{FiniteGluing.glue_restrict} glues exactly agreeing finite tables by
the first applicable patch.  The earlier analytic embedding packet also
constructs simple analytic branches under a derivative separation
certificate.  We reuse these principles.  The additions here are the
completion interface that finishes the value-identification step, a
quantitative continuous gluing theorem with approximate data, and a
singular inverse example that needs a H\"older modulus rather than a
nonzero derivative.

\subsection{Metric presentations and the scope of the statements}

Put $\anE_n=2^{-n}$.  For the exact rational core, a metric presentation
consists of a decidable set $D$ of natural-number codes and a fixed
primitive-recursive rational function $d:D\times D\to\Q_{\ge0}$ satisfying
symmetry, $d(x,x)=0$, and the triangle inequality.  Distinct codes may have
distance zero.  All membership and metric laws have explicitly specified
arithmetic formulas; an application must prove them or retain them as
premises.  We do not derive a metric law from the fact that its checker
terminates.

A regular name is a sequence $x=(x_n)_{n\ge0}$ in $D$ with
\begin{equation}
 d(x_m,x_n)\le\anE_m+\anE_n.
 \label{eq:an-regular}
\end{equation}
For internal use the sequence is a fixed primitive-recursive evaluator
$a(p,n)$, where $p$ is an ordinary numerical parameter.  Quantifying over
all names is external notation.  Each arithmetic theorem is a schema in
the fixed evaluator and its supplied contracts; it does not evaluate an
arbitrary code for a primitive-recursive program.

The semantic statements below also hold for arbitrary metric spaces.
For a general effective presentation, the distance itself has rational
approximants $\Delta(x,y,t)$ with error at most $\anE_t$.  Then
$d(x,y)\le q$ abbreviates
\[
 \forall t\quad \Delta(x,y,t)\le q+\anE_t.
\]
The exact rational version makes the arithmetic kernel easiest to see.
The extension to distance names is addressed after the main construction.
No induction is performed on a statement quantifying over real points.

\begin{lemma}[Removing a vanishing rational error]
\label{lem:an-small-error}
For rational numbers $u,v$, if $u\le v+C\anE_n$ for every natural $n$,
where $C\ge0$ is rational, then $u\le v$.  The contradiction index can be
chosen primitive recursively from the rational codes of $u,v,C$.
\end{lemma}

\begin{proof}
Rational order is decidable.  In the case $u>v$, write $u-v=a/b$ with
positive integers $a,b$, so $u-v\ge1/b$.  Choose an integer $K\ge C$ and
put $n=Kb+1$.  The elementary inequality $2^n> Kb$ gives
$C\anE_n<1/b$ when $K>0$, contradicting the assumed inequality.  If
$K=0$, the contradiction is immediate.  Integer arithmetic, bounded
rounding, and the proof of $2^n\ge n+1$ use quantifier-free induction.
The universal premise is instantiated once, at the displayed numerical
index; it is not the induction formula.
\end{proof}

Define the represented distance bound
\begin{equation}
 \anH d(x,y)\le q
 \quad\Longleftrightarrow\quad
 \forall m,n\;
 d(x_m,y_n)\le q+\anE_m+\anE_n,
 \qquad q\in\Q_{\ge0}.
 \label{eq:an-completion-order}
\end{equation}
Write $x\sim y$ when this holds with $q=0$.  Reflexivity and symmetry
follow immediately.  For transitivity, insert an intermediate sample
$y_t$ and use the two hypotheses to obtain
$d(x_m,z_n)\le\anE_m+\anE_n+2\anE_t$; then apply
Lemma~\ref{lem:an-small-error}.  The same proof gives the triangle
inequality for the represented distance bounds.  Thus regular names
modulo $\sim$ form the represented completion $\anH D$.  Here
completeness is asserted for Cauchy families with a supplied joint family
of names, as in Theorem~\ref{thm:an-diagonal}.  No choice of representatives
for an arbitrary abstract sequence of quotient points is inferred.
Its distance is
represented by the rational samples
\[
 d(x_{n+2},y_{n+2}).
\]
The reverse triangle inequality bounds their pairwise difference by
$2\anE_{m+2}+2\anE_{n+2}\le\anE_m+\anE_n$.
Constant names embed $D$ into $\anH D$, and
\begin{equation}
 \anH d(x,\operatorname{const}(x_n))\le\anE_n.
 \label{eq:an-density}
\end{equation}
Indeed, the defining inequalities are precisely regularity, with an
additional nonnegative error allowance for the constant name.

\begin{theorem}[Explicit diagonal completion]
\label{thm:an-diagonal}
Let $x_{i,p}\in D$ be a fixed two-parameter family such that every
row $p\mapsto x_{i,p}$ is a regular name, and suppose
\[
 d(x_{i,p},x_{j,q})
 \le\anE_i+\anE_j+\anE_p+\anE_q
 \quad\text{for all }i,j,p,q.
\]
The diagonal $z_n=x_{n+2,n+2}$ is a regular name, and its completed
distance from the $j$th name is at most $\anE_j$.  In particular, a
regular sequence of regular names with this explicit joint evaluator
can be flattened primitive recursively.
\end{theorem}

\begin{proof}
Substitute $(i,p)=(m+2,m+2)$ and $(j,q)=(n+2,n+2)$ to obtain
$d(z_m,z_n)\le2\anE_{m+2}+2\anE_{n+2}$, which is at most the
regularity budget.  For the distance to the $j$th name, substitute only
$(i,p)=(m+2,m+2)$:
$d(z_m,x_{j,q})\le2\anE_{m+2}+\anE_j+\anE_q
\le\anE_m+\anE_j+\anE_q$.
This is exactly~\eqref{eq:an-completion-order}.  There is one fixed
joint evaluator in the hypothesis; the proof never chooses a name for
each member of an arbitrary sequence of completed points.
\end{proof}

\subsection{Extension, continuity, and uniqueness}

Let $f:D\to E$ be a fixed primitive-recursive map between rational
metric presentations.  A dyadic uniform-continuity modulus is a function
$\mu:\N\to\N$ satisfying
\begin{equation}
 d_D(u,v)\le\anE_{\mu(k)}
 \quad\Longrightarrow\quad
 d_E(f(u),f(v))\le\anE_k.
 \label{eq:an-map-modulus}
\end{equation}
Replace $\mu(k)$ by $\max_{j\le k}\mu(j)$ to make it nondecreasing.
This replacement is primitive recursive, and its maximum invariant is
quantifier free.  Define
\begin{equation}
 h(n)=\mu(n+2)+2,
 \qquad (\mathsf{Ext}_f x)_n=f(x_{h(n)}).
 \label{eq:an-extension}
\end{equation}

\begin{theorem}[Quantitative universal extension]
\label{thm:an-extension}
Formula~\eqref{eq:an-extension} produces a regular name, respects
$\sim$, and extends $f$ on constant names.  The induced map
$\anH f:\anH D\to\anH E$ has modulus
\[
 \nu(k)=\mu(k+2)+2.
\]
It is the unique uniformly continuous extension of $f$ to these
completions.  All sample constructions and index changes are primitive
recursive when the input routines and moduli are fixed primitive
recursive functions.
\end{theorem}

\begin{proof}
If $m\ge n$, monotonicity of $h$ and regularity give
\[
 d_D(x_{h(m)},x_{h(n)})
 \le2\anE_{h(n)}
 =\tfrac12\anE_{\mu(n+2)}.
\]
The modulus therefore bounds the output distance by $\anE_{n+2}$,
which is stronger than the required regularity inequality.

Suppose now $\anH d_D(x,y)\le\anE_{\nu(k)}$.  For $m,n\ge k$,
the corresponding input distance is at most
\[
 \anE_{\nu(k)}+\anE_{h(m)}+\anE_{h(n)}
 \le 3\anE_{\mu(k+2)+2}
 \le\anE_{\mu(k+2)}.
\]
Consequently the distance between these output samples is at most
$\anE_{k+2}$.  Insert such sufficiently late samples between any two
fixed output samples, use their regularity, and let their indices tend
to infinity by Lemma~\ref{lem:an-small-error}.  This proves that the
completed output distance is at most $\anE_{k+2}$, hence at most
$\anE_k$.  If $x\sim y$, the same argument applies for every $k$ and
proves output equivalence.  Constant names are preserved by inspection.

For uniqueness, let $g$ be another extension with a supplied uniform
modulus $\lambda$.  Fix a name $x$ and an output precision $k$.
Choose
$n\ge\max(\nu(k+2),\lambda(k+2))$.  By
\eqref{eq:an-density}, both $\anH f(x)$ and $g(x)$ are within
$\anE_{k+2}$ of the common value $f(x_n)$.  Their mutual distance is
at most $2\anE_{k+2}$ for every $k$, so they are equivalent.
This argument uses a maximum of two known indices, not a search for a
continuity modulus.
\end{proof}

Identity maps and composition are preserved by this extension, as
equalities of represented maps.  For composition, both candidate maps
are uniformly continuous and agree on every constant name; uniqueness
applies.  Equality of the raw output programs is neither needed nor
generally true.  This is the useful categorical universal property of
completion at the interface of the arithmetic and semantic layers.

\paragraph{Distance names.}
The proof used only triangle inequalities, uniform-continuity
implications, and explicit positive slack.  With named distances,
replace every final bound $d\le q$ by all rational tests
$\Delta_t\le q+\anE_t$.  When an intermediate inequality contains a
sum of $r$ distance observations, evaluate each to a precision $t$ for
which $r\anE_t$ is below the allocated slack.  The displayed factors
$1/2$ and $3/4$ provide such slack; non-strict limiting bounds are
closed with Lemma~\ref{lem:an-small-error}.  The query index is computed
by bounded doubling from the number of terms and requested precision.
For real-valued semantic distances, this lemma is applied to the
rational distance tests, not directly to an unencoded real quantity.
The triangle and modulus laws for $\Delta$ are still explicit
hypotheses.  This gives the same extension schema for effective metric
presentations, including function spaces whose sup distances have only
certified rational approximants.

\section{Quantitative gluing of continuous local data}
\label{sec:analytic-gluing}

Exact membership in a chart is often a poor interface for a represented
point near the chart boundary.  For a convex target, a finite weighted
construction removes the need to select one chart.  Unlike exact
finite-table gluing, it also remains quantitatively stable when the
local values are only approximately compatible.

\subsection{Data and the rational algorithm}

Let $D$ be a rational metric presentation, let $c_1,\ldots,c_N\in D$,
and let $r_i>0$ be rational.  Define
\[
 w_i(x)=\max(0,r_i-d(x,c_i)),
 \qquad S(x)=\sum_{i=1}^N w_i(x).
\]
Fix a presented subdomain $D_0\subseteq D$ and assume a rational
$\delta>0$ satisfies $S(x)\ge\delta$ for every valid code in $D_0$.
All local contracts below are restricted to $D_0$, and every regular
input name used in this section has all its samples in $D_0$.
Thus the completed domain is $\widehat{D_0}$.  Membership of the limit
alone would not justify denominator bounds for the sampled codes.
The implementation retains the ambient metric presentation and a
separate predicate for $D_0$, so centers may lie outside $D_0$ while
remaining valid ambient points.  For each $i$ there is a fixed rational-vector
evaluator $f_i(p,x)\in\Q^d$.  The finite family is either fixed
externally or given by one fixed indexed evaluator; indices do not
encode arbitrary programs.  It is required to be correct only when
$w_i(x)>0$.  Use the $\ell^1$ norm
$\norm{v}_1=\sum_{j<d}|v_j|$, which is an exact rational finite sum.
Assume the finite local contracts
\begin{align}
 \norm{f_i(p,x)}_1&\le M,
 && w_i(x)>0,
 \label{eq:an-local-bound}\\
 \norm{f_i(p,x)-f_i(q,y)}_1
 &\le Ld(x,y)+\anE_p+\anE_q,
 && w_i(x),w_i(y)>0,
 \label{eq:an-local-lip}
\end{align}
where $M,L\ge0$ are rational.  These formulas include precision
coherence by taking $x=y$.  The targets can be real vectors, or complex
numbers represented by pairs of rationals; no complex order is used.

Set a summand to zero when its weight is zero, and otherwise evaluate
the local routine.  Define
\begin{equation}
 A_p(x)=\sum_i w_i(x)f_i(p,x),
 \qquad U_p(x)=\frac{A_p(x)}{S(x)}.
 \label{eq:an-glue}
\end{equation}
This is a total rational algorithm: division has a harmless zero
default outside the certified domain, and the positive denominator
contract proves that the default is never used on that domain.  No
local routine is required to satisfy an analytic assertion outside its
guard.  With exact rational distances the weight comparison is simply
a decidable rational comparison.

\begin{theorem}[Weighted gluing with sample errors]
\label{thm:an-glue}
The algorithm satisfies
\begin{align}
 \norm{U_p(x)}_1&\le M,\label{eq:an-glue-bound}\\
 \norm{U_p(x)-U_q(y)}_1
 &\le C d(x,y)+\anE_p+\anE_q,
 \qquad C=L+\frac{2NM}{\delta}.
 \label{eq:an-glue-lip}
\end{align}
If $j$ is active at $x$ and every active $i$ satisfies
\begin{equation}
 \norm{f_i(p,x)-f_j(q,x)}_1
 \le\eta+\anE_p+\anE_q,
 \label{eq:an-overlap}
\end{equation}
then
\begin{equation}
 \norm{U_p(x)-f_j(q,x)}_1
 \le\eta+\anE_p+\anE_q.
 \label{eq:an-glue-restrict}
\end{equation}
The error does not acquire a factor $N$ in this last estimate.
\end{theorem}

\begin{proof}
The reverse triangle inequality and the fact that the positive-part
function is $1$-Lipschitz give
$|w_i(x)-w_i(y)|\le d(x,y)$ and
$|S(x)-S(y)|\le Nd(x,y)$.  The coefficients $w_i(x)/S(x)$ are
nonnegative and sum to one, proving~\eqref{eq:an-glue-bound}.

Write $\alpha=w_i(x)$ and $\beta=w_i(y)$.  If both are positive,
split the difference as
\[
 \alpha f_i(p,x)-\beta f_i(q,y)
 =(\alpha-\beta)f_i(q,y)
   +\alpha\bigl(f_i(p,x)-f_i(q,y)\bigr).
\]
Its norm is at most
$M|\alpha-\beta|+
\alpha(Ld(x,y)+\anE_p+\anE_q)$.
If $\alpha=0$ or $\beta=0$, the same bound remains valid: the
nonzero summand has norm at most $M|\alpha-\beta|$, and the extra
term is nonnegative.  Thus no assertion about an inactive local
value has been used.  Summing gives
\[
 \norm{A_p(x)-A_q(y)}_1
 \le NMd(x,y)+S(x)
       \bigl(Ld(x,y)+\anE_p+\anE_q\bigr).
\]
Finally,
\[
 \frac{A_p(x)}{S(x)}-\frac{A_q(y)}{S(y)}
 =\frac{A_p(x)-A_q(y)}{S(x)}
  +A_q(y)\frac{S(y)-S(x)}{S(x)S(y)}.
\]
Since $\norm{A_q(y)}_1\le MS(y)$ and $S(x)\ge\delta$, the first
term contributes at most $(L+NM/\delta)d(x,y)+\anE_p+\anE_q$,
and the second contributes at most $(NM/\delta)d(x,y)$.
This proves~\eqref{eq:an-glue-lip}.
To prove~\eqref{eq:an-glue-restrict}, subtract $f_j(q,x)$ inside
the convex sum and apply~\eqref{eq:an-overlap} term by term.
\end{proof}

\subsection{Completion, coverage certificates, and exact gluing}

Choose a natural $c\ge2$ with $C\le2^{c-2}$.  For a regular input
name $x$, define the output samples
\begin{equation}
 b_n=U_{n+2}(x_{n+c}).
 \label{eq:an-glue-name}
\end{equation}
Equation~\eqref{eq:an-glue-lip} gives
\[
 \norm{b_m-b_n}_1
 \le C2^{-c}(\anE_m+\anE_n)
      +\anE_{m+2}+\anE_{n+2}
 \le\tfrac12(\anE_m+\anE_n).
\]
Hence the output is a regular name.  Applying the same finite estimate
to equivalent input names proves extensionality.  The induced map on
the completion is $C$-Lipschitz; the proof follows by inserting late
samples and removing the vanishing error.  This argument directly
allows approximate local evaluators, so it does not require an exact
rational value of the local functions.

The denominator hypothesis can itself be certified by a finite net.
Suppose a finite list $v_1,\ldots,v_s$ covers the domain within distance
$\rho$, and a finite calculation verifies
\[
 S(v_j)\ge\delta+N\rho\quad(j<s).
\]
For any covered point choose its supplied covering witness, use the
$N$-Lipschitz bound for $S$, and obtain $S(x)\ge\delta$.  Internally
the coverage witness is a bounded search result or part of the domain
contract.  A claimed arbitrary open cover does not manufacture this
net or its coverage proof.

If the fixed local routines agree on every overlap to all requested
precisions, use~\eqref{eq:an-glue-restrict} with
$\eta=\anE_k$ for arbitrary $k$.  The completed global value equals
each active local value.  At a completed point with positive clearance
inside a chart, sufficiently late rational samples remain in that
chart, so the same equality holds there.  Positive clearance follows
from the robust cover data for at least one chart; exact chart identity
need not be selected by the program.

If one has only a fixed overlap tolerance $\eta$, the theorem gives
approximate gluing with exactly that tolerance.  A different family of
local answers at every precision must additionally have a cross-stage
coherence contract.  Same-stage agreement is insufficient: all charts
could agree on $0$ at even stages and on $1$ at odd stages.  Thus the
construction does not turn independent approximate local solutions
into an exact global solution without a coherence premise.

\paragraph{Named distances and safe support tests.}
If $d(x,c_i)$ is approximated by $\Delta_i(t)$ with error
$\anE_t$, use the inner weight
\[
 \widetilde w_i=\max(0,r_i-\Delta_i(t)-\anE_t).
\]
Then $0\le\widetilde w_i\le w_i$,
$w_i-\widetilde w_i\le2\anE_t$, and
$\widetilde w_i>0$ certifies genuine membership in the local domain.
If $2N\anE_t\le\delta/2$, their sum is at least $\delta/2$.
Replacing the exact weights by these inner weights changes the
normalized sum by at most
\[
 \frac{8MN}{\delta}\anE_t.
\]
Indeed, the numerator changes by at most
$M\sum_i(w_i-\widetilde w_i)$; the denominator correction contributes
the same amount before division by a denominator at least $\delta/2$.
Choose $t$ so this additional error fits the requested output budget.
This supplies a finite safe algorithm even when exact chart membership
is not decidable from a metric name.

The target convexity is essential.  Weighted averages of manifold
points need not lie on the manifold.  A manifold-valued version needs
an explicit retraction from the convex neighborhood, together with its
domain margin and continuity modulus.  None is inferred from the
word ``manifold.''

\section{From approximate solutions to an identified completed point}
\label{sec:residual-recovery}

The repository's coherence obstruction prevents an unrestricted
inference from finite solvability at every horizon to one computable
coherent branch.  A positive theorem must specify the extra information
that rules out that obstruction.  Here it is a modulus bounding the
diameter of the approximate solution sets.  This is the same core
principle already used by \texttt{CoherentMinimizer}; completion now
also identifies the limiting value and proves its uniqueness.

\subsection{The generic residual theorem}

Let $r:D\to\Q_{\ge0}$ be a fixed uniformly continuous residual, with
a supplied dyadic modulus.  Assume a fixed primitive-recursive
procedure $a:\N\to D$ and a primitive-recursive function $\kappa$
satisfy
\begin{align}
 r(a(n))&\le\anE_n,\label{eq:an-residual-witness}\\
 r(u),r(v)\le\anE_{\kappa(k)}
 &\quad\Longrightarrow\quad d(u,v)\le\anE_k.
 \label{eq:an-diameter}
\end{align}
The approximate witnesses are an assumption to prove in the application;
they are not extracted from an assertion that an exact zero exists.
Normalize $\kappa$ to the nondecreasing majorant
\[
 \kappa^*(k)=k+\max_{j\le k}\kappa(j).
\]
It is at least $k$ and at least $\kappa(k)$.

\begin{theorem}[Recovery and identification by a diameter modulus]
\label{thm:an-residual-recovery}
The samples
\[
 z_n=a(\kappa^*(n+2))
\]
form a regular name.  The continuous extension $\anH r$ satisfies
$\anH r(z)=0$.  This is the unique completed zero of $\anH r$.
Every construction is primitive recursive relative to the fixed
displayed procedures.  The same statement holds for a named residual
when the finite witness and diameter contracts include its sample
error allowances.
\end{theorem}

\begin{proof}
For $m\ge n$, both residuals are at most
$\anE_{\kappa^*(n+2)}\le\anE_{\kappa(n+2)}$.
Equation~\eqref{eq:an-diameter} therefore gives
$d(z_m,z_n)\le\anE_{n+2}$, proving regularity.
Also $r(z_n)\le\anE_{n+2}$ because $\kappa^*(n+2)\ge n+2$.

Let $\nu_r$ be the extension modulus furnished by
Theorem~\ref{thm:an-extension}.  To bound $\anH r(z)$ at output
precision $k$, choose
$n\ge\max(k+2,\nu_r(k+2))$.  Density gives
$\anH d(z,\operatorname{const}(z_n))\le\anE_n$, so
$|\anH r(z)-r(z_n)|\le\anE_{k+2}$.
Together with $r(z_n)\le\anE_{k+2}$ this bounds
$|\anH r(z)|$ by $2\anE_{k+2}$ for every $k$.  Hence it is zero.
This is the canonical-evaluation step missing from a Cauchy estimate
alone.

For uniqueness, let $y$ be a regular name with $\anH r(y)=0$.
Fix $k$ and put $t=\kappa^*(k+3)$.  Choose
\[
 n\ge\max(k+3,\nu_r(t+1)),\qquad
 m\ge\max(k+3,t).
\]
Uniform continuity and density imply
$r(y_n)\le\anE_{t+1}$, while
$r(z_m)\le\anE_{m+2}\le\anE_t$.
Both are below $\anE_{\kappa(k+3)}$, so
$d(z_m,y_n)\le\anE_{k+3}$.
The completed distance between $z$ and $y$ is consequently at most
\[
 \anE_m+\anE_{k+3}+\anE_n
 \le3\anE_{k+3}<\anE_k.
\]
This holds at every $k$ and proves $y\sim z$.  Every index in the
proof is an explicitly displayed primitive-recursive composition or
maximum.  No exact solution was chosen at an intermediate stage.
\end{proof}

There are two separate assumptions here.  Approximate solvability
supplies points; the diameter modulus forces any sufficiently accurate
points to agree.  Neither replaces the other.  Mere uniqueness of a
semantic zero does not automatically supply a primitive-recursive
diameter modulus, and a diameter bound without approximate solvability
could be vacuous.

\subsection{A complete contraction instance}

Let $F:D\to D$ be a fixed map with
$d(Fu,Fv)\le qd(u,v)$ for a supplied rational $0\le q<1$.
Choose a starting point $x_0$ and integer $B\ge0$ such that
$d(x_0,Fx_0)\le2^B$.  Let $H\ge1$ be an integer satisfying
$H(1-q)\ge1$.  Both $B$ and $H$ can be found by bounded searches from
the rational data.

For $\gamma=1-q$, the identity
$q(1+\gamma)=1-\gamma^2\le1$ and the finite Bernoulli inequality
$(1+\gamma)^H\ge1+H\gamma\ge2$ imply $q^H\le1/2$.
The Bernoulli inequality is proved by induction on the exponent with
a polynomial inequality as invariant.  Iterating the contraction
gives
\[
 d(F^j x_0,F^{j+1}x_0)\le q^j2^B.
\]
For the residual $r(x)=d(x,Fx)$, set
\[
 a(n)=F^{H(n+B)}x_0.
\]
Then $r(a(n))\le\anE_n$.  This constructs approximate solutions
directly from finite iteration; it assumes no fixed point.

The triangle inequality gives
\[
 (1-q)d(u,v)\le r(u)+r(v).
\]
Choose $c$ with $2^{1-c}\le1-q$.  Then
$\kappa(k)=k+c$ is a diameter modulus.
Moreover $|r(u)-r(v)|\le(1+q)d(u,v)$, so the residual is uniformly
continuous.  Theorem~\ref{thm:an-residual-recovery} now constructs a
unique completed zero.  By the extension theorem, the residual
identifies with $x\mapsto\anH d(x,\anH F(x))$: these two continuous
functions agree on the dense constant names.  Thus the zero is the
unique fixed point of $\anH F$.

The source module \texttt{ContractionIterates} already supplies
quantifier-free induction proofs of contraction estimates for
iterations.  The additional reusable step here is the completed
fixed-point identity and its uniqueness, with the exact dependency
on continuity and density exposed.

\section{A singular inverse that does not need derivative separation}
\label{sec:singular-cube-root}

Derivative separation is a sufficient way to obtain a diameter
modulus, but it is not necessary.  The real cube root is continuous at
zero even though the derivative of $y\mapsto y^3$ vanishes there and
the inverse is not locally Lipschitz.  Its quantitative inverse
modulus is still primitive recursive.

The inspected modules \texttt{StableCubicRoot} and
\texttt{CubicInverse} concern $y^3+y=x$, whose extra linear term
gives inverse Lipschitz constant one.  The example here deliberately
uses $y^3=x$ and includes negative inputs and the singular value zero.
It does not produce a holomorphic branch of a complex cube root near
zero, nor does it decide the sign of a represented real number.

\begin{lemma}[A global cubic separation inequality]
\label{lem:an-cubic-separation}
For all rational $u,v$,
\begin{equation}
 |u-v|^3\le4|u^3-v^3|.
 \label{eq:an-cubic-separation}
\end{equation}
The same inequality holds for represented reals by approximation.
\end{lemma}

\begin{proof}
If $u,v$ have the same sign, change both signs if necessary and assume
$u\ge v\ge0$.  Then
\[
 u^3-v^3=(u-v)(u^2+uv+v^2)\ge(u-v)^3,
\]
because $u^2+uv+v^2-(u-v)^2=3uv\ge0$.
If they have opposite signs, put $s=|u|$ and $t=|v|$.  The exact
identity
\[
 4(s^3+t^3)-(s+t)^3=3(s+t)(s-t)^2\ge0
\]
proves the result.  All case splits compare rational integers, and
all substantive steps are polynomial identities and nonnegative
products.  Passing to real names uses the explicit product error
bounds on a supplied bounded interval and
Lemma~\ref{lem:an-small-error}.
\end{proof}

Let $A_p$ be a fixed regular name of $a\in[-1,1]$.  Semantically
$|A_p-a|\le\anE_p$; internally this is expressed by the
corresponding cross-sample inequalities.  Define
$C_p=\max(-1,\min(1,A_p))$.  Clipping is $1$-Lipschitz and fixes
$[-1,1]$, so $|C_p-a|\le\anE_p$.

For a requested residual index $n$, put
\[
 p=n+4,\qquad J=2^{n+5},\qquad q_j=-1+2j/J\quad(0\le j\le J).
\]
Use bounded search to select the largest $j\le J$ for which
$q_j^3\le C_p$ and return $a_n=q_j$.  There is always at least one
candidate because $q_0=-1$ and $C_p\ge-1$.  All comparisons concern
the rational sample $C_p$, so even the input $a=0$ causes no exact
real comparison.

\begin{proposition}[Primitive-recursive approximate cubic roots]
\label{prop:an-cubic-approx}
The returned point lies in $[-1,1]$ and satisfies
\begin{equation}
 |a_n^3-a|\le\anE_{n+2}.
 \label{eq:an-cubic-residual}
\end{equation}
\end{proposition}

\begin{proof}
If $j=J$, maximality forces $C_p=1$ and the error relative to
$C_p$ is zero.  Otherwise
$q_j^3\le C_p<q_{j+1}^3$.  For $x,y\in[-1,1]$,
\[
 |x^3-y^3|=|x-y||x^2+xy+y^2|\le3|x-y|.
\]
The mesh length is $2/J$, so
$0\le C_p-a_n^3\le6/J$.
Adding the input sample error yields
\[
 |a_n^3-a|\le6\,2^{-(n+5)}+2^{-(n+4)}
 =\tfrac14\,2^{-n}=\anE_{n+2}.
\]
The maximum-search invariant states that the selected index lies
within the finite grid, passes the comparison, and every later
inspected index fails it.  A Boolean fold expresses the last condition;
the invariant is quantifier free.  No intermediate-value theorem or
compactness principle has been used to produce the witness.
\end{proof}

If $|u^3-a|$ and $|v^3-a|$ are at most
$\anE_{3k+3}$, then
\[
 |u-v|^3\le4\bigl(|u^3-a|+|v^3-a|\bigr)
 \le8\anE_{3k+3}=\anE_k^3.
\]
Monotonicity of the cube on nonnegative rationals proves
$|u-v|\le\anE_k$.  Thus
\begin{equation}
 \kappa(k)=3k+3
 \label{eq:an-cubic-kappa}
\end{equation}
is an explicit residual-diameter modulus.  The finite contractions in
this proof remain valid with rational sample errors: use an arbitrary
sufficiently precise common sample $A_t$, apply the rational triangle
inequality, and remove $2\anE_t$ by
Lemma~\ref{lem:an-small-error}.

In particular, the simple schedule
\[
 z_k=a_{3k+9}
\]
already gives a regular name: for $m\ge k$ both residuals meet the
threshold $\kappa(k+2)=3k+9$, so
$|z_m-z_k|\le\anE_{k+2}$.
The residual is continuous on $[-1,1]$, with Lipschitz constant three,
so the recovery proof identifies $z^3=a$ and proves uniqueness.
For a supplied integer bound $B\ge1$ with $|a|\le B$, apply the
construction to $a/B^3\in[-1,1]$ and multiply the output by $B$.
An extra sample shift $s$ with $B\le2^s$ restores the regularity
budget.  This extends the construction to every bounded real-name
family using only rational scaling and bounded doubling.

Moreover, for the outputs $x=\sqrt[3]{a}$ and $y=\sqrt[3]{b}$,
\[
 |x-y|^3\le4|a-b|.
\]
Thus the inverse itself has dyadic modulus $k\mapsto3k+2$.
There is no discontinuous branch decision.  A loss from $k$ output
bits to roughly $3k$ input bits expresses the conditioning at the
singular value.  This resolves the construction at this singularity
without making the much stronger, false assertion that every
singular analytic branch can be chosen continuously.

\section{Formalization obligations for the continuous layer}
\label{sec:analytic-lean-obligations}

The accompanying Lean blueprint names the semantic interfaces and the
finite statements to implement.  These are not completed new
\texttt{Prf PRAAx} proofs.  The reliable order of implementation is
as follows.

\begin{enumerate}
\item Encode rational vectors, their $\ell^1$ norm, the positive-part
operation, and the finite sums in~\eqref{eq:an-glue}.  Denominators
are positive natural codes; signed numerators are handled by the
existing integer arithmetic layer.  Prove the one-summand inequality
in the proof of Theorem~\ref{thm:an-glue} first.

\item Derive the finite weighted-sum bound by induction on the list
length.  The invariant is the cleared rational inequality for the
partial numerator and partial weight sum.  Retain $\delta>0$ as an
explicit premise before cancelling denominators.  This produces an
object theorem for arbitrary finite chart tables and precision
parameters.

\item Formalize the reindexing inequalities for
Theorem~\ref{thm:an-extension} and the diagonal theorem.  Their
internal conclusions are universal closures of rational inequalities
between samples of specified \texttt{PR} programs.  Regularity and
modulus hypotheses are instantiated at the displayed indices.
There is no induction over names or real numbers.

\item Reuse the maximum-normalization and coherence lemmas from
\texttt{CoherentMinimizer}; add canonical-evaluation and uniqueness
lemmas using the completion interface.  The proof goal is the
explicit sample equality formula for the residual name, not merely a
Lean equality in an ambient complete space.

\item Obtain the two cubic identities through the existing polynomial
derivation mechanism.  Derive the same-sign and opposite-sign cases
using decidable rational sign tests.  Implement the finite grid
selector with the existing bounded-search interface; discharge its
endpoint, maximality, mesh, and residual inequalities.  Finally
compose the residual contracts with the explicit cubic diameter
modulus.

\item Only after these deep arithmetic endpoints are available,
state the external semantic theorems about completions, continuous
maps, and real cube roots.  Prove that the represented evaluations
agree with those semantic objects.  An ambient Lean theorem alone
does not replace any of the preceding endpoints.
\end{enumerate}

For a general distance presentation, every object formula additionally
contains the finite distance-query index and its error allowance.
The trust boundary includes the proofs of those distance contracts,
of coverage, and of the local moduli.  These are useful explicit
obligations, rather than implicit appeals to compactness or a
qualitative continuity theorem.
